% Generated draft: requires mathematical and source review.
\documentclass[11pt]{article}
\usepackage[margin=25mm]{geometry}
\usepackage{fontspec}
\setmainfont{Latin Modern Roman}
\usepackage{amsmath,amssymb,mathtools,unicode-math}
\setmathfont{Latin Modern Math}
\usepackage{enumitem,needspace,xurl,hyperref}
\setlength{\parindent}{0pt}
\setlength{\parskip}{5pt}
\setlength{\emergencystretch}{3em}
\newcommand{\N}{\mathbb{N}}
\newcommand{\Z}{\mathbb{Z}}
\newcommand{\Q}{\mathbb{Q}}
\newcommand{\R}{\mathbb{R}}
\newcommand{\C}{\mathbb{C}}
\newcommand{\F}{\mathbb{F}}
\newcommand{\A}{\mathbb{A}}
\newcommand{\PP}{\mathbb{P}}
\newcommand{\E}{\mathbb{E}}
\newcommand{\eps}{\varepsilon}
\newcommand{\dd}{\,\mathrm{d}}
\begin{document}
\section*{UnsolvedMath ID: 700787}
\textbf{Source problem code:} catalogue entry
\subsection{Definitions and mathematical statement}
Let \(p\) be a prime. Interpret a \(p\)-group in the usual, possibly infinite, sense: a group in which every element has finite order equal to a power of \(p\). The source does not explicitly specify this convention or impose finiteness on \(G\).

For \(a,g\in G\), write \(a^g=g^{-1}ag\), and define
\[
 a^G=\{a^g:g\in G\},\qquad
 C_G(a)=\{x\in G:xa=ax\},\qquad
 Z(G)=\{z\in G:zx=xz\text{ for every }x\in G\}.
\]
Here \(\langle a,a^g\rangle\) denotes the subgroup generated by \(a\) and \(a^g\).

Is it true that, for every prime \(p\) and every \(p\)-group \(G\), if there exists \(a\in G\) such that
\[
 \operatorname{ord}(a)=p,\qquad
 \bigl|\langle a,a^g\rangle\bigr|<\infty\quad\text{for every }g\in G,
 \qquad
 |C_G(a)\cap a^G|<\infty,
\]
then
\[
 Z(G)\ne\{1\}?
\]
The supplied source states that the answer is affirmative for \(p=2\); this is retained as a source assertion, without independent verification.
\subsection{Short English statement}
Must a possibly infinite \(p\)-group have nontrivial centre if it contains an element \(a\) of order \(p\) such that every subgroup generated by \(a\) and one conjugate of \(a\) is finite, and only finitely many conjugates of \(a\) commute with \(a\)? The source states an affirmative answer for \(2\)-groups.
\subsection{Sources}
\begin{itemize}
\item Supplied UnsolvedMath local registry record: permanent ID \(700787\), source code \texttt{catalogue entry}, title \emph{Kourovka Notebook Problem 11.95}. The registry labels the problem open; this label is supplied metadata, not an independently verified assessment.
\item \emph{The Kourovka Notebook (Unsolved Problems in Group Theory)}, Problem 11.95, originally in the 11th issue (1990), proposed by S. P. Strunkov. The supplied record identifies the 21st issue edition (2026), arXiv:1401.0300v48. Its collection description names E. I. Khukhro and V. D. Mazurov as editors. This bibliographic attribution and the stated affirmative result for \(2\)-groups have not been independently verified.
\end{itemize}
\end{document}
